\documentclass[profile=standard]{ipara-handbook}
\title{B08｜傅里叶展开与正交基：函数表示与正交投影}
\author{}
\date{}
\begin{document}
\maketitle

\begin{quote}
本册回答：\textbf{傅里叶系数为什么一定长成“函数乘正弦/余弦再积分”的样子？它与线性代数中的坐标、正交投影、最小平方误差究竟是什么关系？}
\end{quote}

\section{本册定位：把函数系数读成正交坐标}
B00 Own 内积、正交与投影；线代 Topic01 Own 基与坐标；高数 Topic11 Own 傅里叶级数的周期延拓、系数计算、收敛与恢复。B08 只 Own：
\[
\boxed{
\text{正交函数系}
\xrightarrow{\text{内积}}
\text{投影坐标}
\xrightarrow{\text{有限部分和}}
\text{正交投影近似}
}
\]

它不负责“级数最终是否逐点等于原函数”；那是高数 Topic11 的收敛问题。

\section{有限维原型：为什么正交坐标能被单独抽取}
若 $e_1,\dots,e_n$ 是正交基，且
\[
v=\sum_{k=1}^n c_ke_k,
\]
与 $e_j$ 取内积：
\[
\langle v,e_j\rangle
=\sum_k c_k\langle e_k,e_j\rangle
=c_j\langle e_j,e_j\rangle.
\]
所以
\[
\boxed{c_j=\frac{\langle v,e_j\rangle}{\langle e_j,e_j\rangle}.}
\]

这条公式是 B08 的母原型。核心不是“内积能算系数”，而是\textbf{正交性让其他方向的交叉项全部消失}。

\section{把向量换成函数，点积自然变成积分内积}
在 $[-L,L]$ 上定义
\[
\langle f,g\rangle=\int_{-L}^{L}f(x)g(x)dx.
\]
三角函数满足正交关系，例如对 $m\neq n$，
\[
\int_{-L}^{L}\cos\frac{m\pi x}{L}\cos\frac{n\pi x}{L}dx=0.
\]

于是当我们用
\[
1,\quad \cos\frac{n\pi x}{L},\quad \sin\frac{n\pi x}{L}
\]
作为正交函数系时，系数必然由内积投影产生：
\[
\boxed{
a_n=\frac1L\int_{-L}^{L}f(x)\cos\frac{n\pi x}{L}dx,
}
\]
\[
\boxed{
b_n=\frac1L\int_{-L}^{L}f(x)\sin\frac{n\pi x}{L}dx.
}
\]
常数项对应常数方向的投影。

因此“乘一个三角函数再积分”不是傅里叶的神秘技巧，而是在函数空间里做坐标读取。

\section{有限部分和：不仅是截断，更是最佳平方误差投影}
令 $V_N$ 为前若干个三角基函数张成的有限维子空间，$S_Nf$ 为对应傅里叶部分和。由于
\[
f-S_Nf\perp V_N,
\]
对任意 $p\in V_N$，有
\[
f-p=(f-S_Nf)+(S_Nf-p),
\]
两项正交，所以
\[
\lVert f-p\rVert_2^2
=\lVert f-S_Nf\rVert_2^2+\lVert S_Nf-p\rVert_2^2
\ge \lVert f-S_Nf\rVert_2^2.
\]
于是
\[
\boxed{S_Nf\text{ 是 }V_N\text{ 中对 }f\text{ 的最佳 }L^2\text{ 近似}.}
\]

这一步把 B00 的“投影 = 最近点”完整迁移到函数空间，也是 B08 比单纯系数推导更深的一层。

\section{奇偶性为什么能自动消掉一半系数}
若 $f$ 为偶函数，则
\[
f(x)\sin\frac{n\pi x}{L}
\]
为奇函数，所以
\[
b_n=0.
\]
若 $f$ 为奇函数，则常数方向和余弦方向投影均为零。

因此“偶函数只有余弦项、奇函数只有正弦项”不是一个独立口诀，而是\textbf{对称性使某些内积自动为零}。

做题时真正要检查的是完整被积函数的奇偶性，而不是只看 $f$ 本身。

\section{为什么“算出坐标”不等于“级数已经恢复原函数”}
有限维空间里，给出一组基坐标就能精确恢复向量；无限维函数空间没有这么简单。

写出傅里叶系数只说明我们得到了每个正交方向上的投影坐标。要从这些坐标重建 $f$，还要处理：
\begin{itemize}
\item 所选正交系统是否足够完备；
\item 级数在哪种意义下收敛；
\item 逐点收敛时连续点、跳跃点分别恢复什么值；
\item 是否允许逐项积分、逐项求导。
\end{itemize}

所以必须阻断：
\[
\boxed{\text{系数存在}\neq\text{逐点恢复成立}.}
\]

高数 Topic11 Own 这些收敛资格。

\section{从 Bessel 到 Parseval：投影坐标还保存多少能量}
对任意有限个规范正交方向，正交投影满足
\[
\sum_{k=1}^N|c_k|^2\le \lVert f\rVert_2^2.
\]
这就是 Bessel 型不等式的几何含义：投影到若干方向得到的平方坐标总量，不会超过原对象的平方范数。

若正交系统完备并在相应 Hilbert 空间意义下展开完整，则进一步有 Parseval 等式
\[
\lVert f\rVert_2^2=\sum_k|c_k|^2.
\]

数学一不需要完整 Hilbert 空间理论，但这条 Extension 能解释：傅里叶系数不是一串任意常数，而是函数“能量”在正交方向上的分配。

\section{母例：偶延拓为什么只需余弦投影}
在 $[0,L]$ 上给定函数 $f$，做偶延拓到 $[-L,L]$。偶延拓后的函数与所有正弦基正交，因此
\[
b_n=0.
\]
余弦系数为
\[
a_n=\frac2L\int_0^Lf(x)\cos\frac{n\pi x}{L}dx.
\]
系数前的 $2/L$ 并不是需要死记的新常数，而来自：
\begin{itemize}
\item 全区间内积利用偶性缩成两倍半区间积分；
\item 再除以余弦基函数在当前内积下的范数平方。
\end{itemize}

因此若区间或归一化约定改变，应重新从投影公式生成，而不是硬套原系数。

\section{B08 串起哪些章节}
\begin{tabular}{p{0.23\linewidth}p{0.30\linewidth}p{0.39\linewidth}}
\hline
Owner / 上游 & 提供对象 & B08 的翻译责任\\
\hline
B00 / 线代 Topic01 & 内积、正交基、投影坐标、最近点 & 把有限维投影机制迁移到函数对象\\
高数 Topic11 & 傅里叶基、延拓、系数、收敛 & 解释系数来源与有限部分和的最佳平方近似\\
高数 Topic05 & 积分 & 积分在这里承担“函数内积”而非普通面积累积的角色\\
\hline
\end{tabular}

\section{反桥梁}
\begin{tabular}{p{0.20\linewidth}p{0.04\linewidth}p{0.20\linewidth}p{0.48\linewidth}}
\hline
A & & B & 真正区别\\
\hline
正交函数系 & $\neq$ & 已证明完备 & 正交只消交叉项；完备性决定能否完整恢复\\
傅里叶系数 & $\neq$ & 函数点值 & 系数是全区间投影坐标，不是某点采样\\
有限部分和最佳 $L^2$ 近似 & $\neq$ & 每一点都最接近 & $L^2$ 最优是整体平方误差意义，不是逐点误差最小\\
正交 & $\neq$ & 概率独立 & 一个是函数内积结构，一个是联合概率分解\\
\hline
\end{tabular}

\section{调用信号与一页压缩}
看到“傅里叶系数来源、奇偶项消失、正弦/余弦展开、部分和、特殊点求和”时，先问：
\begin{enumerate}
\item 当前区间与内积是什么？
\item 基函数是否正交，范数平方是多少？
\item 哪些系数因对称性内积为零？
\item 当前任务只需投影系数，还是还要回 Topic11 审计收敛与跳跃点恢复？
\end{enumerate}

最终只需记住
\[
\boxed{
\text{傅里叶系数 = 函数在正交基方向上的投影坐标；部分和 = 对有限基子空间的正交投影。}
}
\]

\end{document}
